%!TEX root = ../main.tex

\chapter{Το Σύστημα torcpy}   

\section{Τι είναι το torcpy}

Πέρα από την αναγκή για λύσεις για την πλήρη αξιοποίησει των ετερογενών πόρων ενός υπολογιστικού κόμβου σε εφαρμογές \textbf{HPC} (π.χ. StarPU), υπάρχει και η ανάγκη για λύσεις που μεγιστοποιούν την απόδοση σε μία συστάδα (cluster) υπολογιστικών κόμβων. Μια προσέγγιση για την επίτευξη αυτού του σκοπού είναι η χρήση του \textbf{MPI} (Message Passing Interface) για την επικοινωνία μεταξύ των κόμβων. Ωστόσο, η χρήση του \textbf{MPI} μπορεί να είναι περίπλοκη και απαιτεί σημαντική προσπάθεια για την ανάπτυξη και τη συντήρηση των εφαρμογών.

Ταυτόχρονα, η συνεχής και ταχεία εξέλιξη και προσθήκη νέων βιβλιοθήκων κάνουν την γλώσσα \textbf{Python} πιο αποτελεσματική για την ανάπτυξη προηγμένων αλγορίθμων και επιστημονικών εφαρμογών. Επιπλέον, προσφέρει διεπαφές υψηλού επιπέδου που διασυνδέουν τους χρήστες με ταχύτατες βιβλιοθήκες HPC γραμμένες σε γλώσσες συστήματος όπως η C, η C++ και η Fortran.

Μία λύση που προσπαθεί να συνδυάσει τα πλεονεκτήματα της \textbf{Python} με την απόδοση του \textbf{MPI} είναι το \textbf{torcpy}. Το torcpy δεν σχεδιάστηκε από το μηδέν ως μια απομονωμένη ιδέα, αλλά αποτελεί τη φυσική εξέλιξη και την αμιγή υλοποίηση σε Python (pure Python implementation) της ισχυρής βιβλιοθήκης χρόνου εκτέλεσης TORC (Task Orchestration and Routing Code), η οποία ήταν γραμμένη σε C/C++ \cite{hadjidoukas2020torcpy}.

Το \textbf{torcpy} είναι ένα σύστημα χρόνου εκτέλεσης ανοιχτού κώδικα, αγνωστικιστικού ως προς την πλατφόρμα (platform-agnostic), το οποίο υλοποιεί προσαρμοστική εξισορρόπηση φόρτου (adaptive load balancing) και ενορχηστρώνει τον χρονοπρογραμματισμό εργασιών τόσο σε συστήματα κοινόχρηστης όσο και κατανεμημένης μνήμης \cite{hadjidoukas2020torcpy}.

Το παρόν κεφάλαιο παρουσιάζει αναλυτικά το σύστημα \textbf{torcpy}, εστιάζοντας στις βασικές του αρχές, την αρχιτεκτονική του και τις δυνατότητες που προσφέρει για την επιτέυξη των στόχων που θέσαμε ως το βασικό κίνητρο για την αναπτυξή του.

\section{Αρχιτεκτονική του torcpy}

Για την επίτευξη των προαναφερθέντων στόχων, το torcpy υλοποιήθηκε ως μια λεπτή αλλά 
εξαιρετικά αποδοτική στιβάδα λογισμικού πάνω από την τεχνολογία του MPI (συγκεκριμένα 
αξιοποιώντας το πακέτο \texttt{mpi4py}) και την τεχνολογία του multithreading της Python 
(αξιοποιώντας τα πακέτα threading και queue).

Η θεμελιώδης αρχιτεκτονική διάταξη μιας εφαρμογής που βασίζεται στο torcpy εδράζεται στην ανάπτυξη πολλαπλών διεργασιών MPI, οι οποίες κατανέμονται στους διαθέσιμους 
υπολογιστικούς κόμβους του συστήματος. Στο εσωτερικό κάθε διεργασίας MPI κατοικεί ένα εξαιρετικά βελτιστοποιημένο υποσύστημα νημάτων (thread subsystem) που αποτελείται από τρεις βασικούς πυλώνες \cite{hadjidoukas2020torcpy}:

\begin{enumerate}
    \item \textbf{Νήματα Εργαστών (Worker Threads):}  Κάθε διεργασία περιλαμβάνει ένα ή περισσότερα worker threads. Ο αποκλειστικός ρόλος τους είναι να εξάγουν (extract) εργασίες από τις διαθέσιμες ουρές και να εκτελούν τον παραγόμενο κώδικα Python. 
    \item \textbf{Νήμα Εξυπηρετητή (Server Thread):} Αποτελεί την καρδιά του συστήματος επικοινωνίας και τη γέφυρα με τον εξωτερικό κόσμο (άλλες διεργασίες MPI). Κάθε 
    διεργασία διαθέτει ένα εξειδικευμένο νήμα διακομιστή (dedicated server thread), το οποίο διαχειρίζεται αποκλειστικά τις απομακρυσμένες λειτουργίες (remote operations) 
    διαχείρισης εργασιών και δεδομένων. Όλα τα μηνύματα που αφορούν επικοινωνία μεταξύ 
    των κόμβων ανταλλάσσονται μέσω αυτού του νήματος με διαφανή προς τον χρήστη τρόπo.
    \item \textbf{Ουρές Εργασιών (Task Queues):} Ένα σύνολο από ιεραρχικές ουρές (set of queues) διατηρείται στην τοπική μνήμη της διεργασίας. Οι εργασίες υποβάλλονται σε αυτές τις ουρές, οι οποίες κατηγοριοποιούνται και οργανώνονται σύμφωνα με το επίπεδο παραλληλισμού (level of parallelism) από το οποίο αρχικά δημιουργήθηκαν. Για παράδειγμα, το αρχικό πρόγραμμα δημιουργεί εργασίες στην ουρά "επιπέδου 1". Αν μια τέτοια εργασία γεννήσει η ίδια νέες εργασίες, αυτές θα τοποθετηθoύν στην ουρά "επιπέδου 2", και ούτω καθεξής.
\end{enumerate}

Στην περίπτωση όπου ο χρήστης εκτελεί την εφαρμογή ως αμιγώς πολυνηματικό 
(multithreaded) κώδικα σε ένα μόνο μηχάνημα —δηλαδή δημιουργώντας μόνο μία διεργασία 
MPI η οποία φιλοξενεί πολλαπλά νήματα εργατών— το torcpy προσαρμόζεται αυτόματα. Σε 
αυτό το σενάριο, όλες ανεξαιρέτως οι λειτουργίες υποβολής, εξαγωγής και μεταφοράς 
δεδομένων πραγματοποιούνται αποκλειστικά μέσω της ταχύτατης τοπικής διαμοιραζόμενης 
μνήμης (shared memory), απαλείφοντας εντελώς το κόστος των κλήσεων του δικτύου MPI \cite{hadjidoukas2020torcpy}.

\subsection{Δομή Εργασιών και Διαχείριση Αποτελεσμάτων}

Μια κρίσιμη απόφαση σχεδιασμού για την εξασφάλιση της ταχύτητας του συστήματος αφορά τη 
δομή των ίδιων των εργασιών. Προκειμένου να αποφευχθεί το μεγάλο υπολογιστικό κόστος 
(overhead) που συνεπάγεται η δημιουργία, η δέσμευση μνήμης και η σειριοποίηση 
(serialization) ογκωδών κλάσεων, οι εργασίες υλοποιούνται αποκλειστικά ως δομές λεξικών της Python (Python dictionaries). Τα λεξικά της Python είναι δομές υψηλής ταχύτητας στο επίπεδο της γλώσσας C (cpython backend), μειώνοντας κατακόρυφα τον χρόνο που απαιτείται για τη δημιουργία της εργασίας \cite{hadjidoukas2020torcpy}.

Όταν η εργασία ολοκληρώνεται, το αποτέλεσμα της συνάρτησης αποθηκεύεται με απόλυτη 
διαφάνεια πίσω στον περιγραφέα της εργασίας (task descriptor ή future). Το εντυπωσιακό είναι ότι το αποτέλεσμα επιστρέφεται και αποθηκεύεται απευθείας στη διεργασία MPI η οποία δημιούργησε (spawned) την εργασία αρχικά, λύνοντας το πρόβλημα της διαχείρισης της μνήμης για τον προγραμματιστή. Ακολουθώντας πιστά το πρότυπο PEP 3148, ο τελικός 
χρήστης μπορεί μετέπειτα να αποκτήσει πρόσβαση στο αποτέλεσμα χρησιμοποιώντας απλώς 
την κλήση της μεθόδου task.result(). Με πανομοιότυπο τρόπο, παρέχεται η δυνατότητα 
πρόσβασης στα αρχικά ορίσματα της εργασίας μέσω της κλήσης task.input() \cite{hadjidoukas2020torcpy}.

\subsection{ Ο Ρόλος και η Λειτουργία του Νήματος Εξυπηρετητή (Server Thread)}

Το Server Thread αποτελεί τον ενορχηστρωτή του συστήματος και λειτουργεί συνεχώς στο 
παρασκήνιο. Η λειτουργία του είναι ασύγχρονη και είναι επιφορτισμένο με τρεις κύριες, ζωτικής σημασίας ευθύνες \cite{hadjidoukas2020torcpy}:

\begin{enumerate}
    \item \textbf{Λήψη και Ένταξη:} Αναλαμβάνει τη λήψη νέων, εισερχόμενων εργασιών που καταφθάνουν από άλλους κόμβους (μέσω του δικτύου MPI) και αναλαμβάνει την τοποθέτησή τους (insertion) στις κατάλληλες τοπικές ουρές της διεργασίας του.
    \item \textbf{Συλλογή Αποτελεσμάτων:} Παραλαμβάνει τις εργασίες που έχουν πλέον ολοκληρωθεί, μαζί με τα αντίστοιχα αποτελέσματά τους, δρομολογώντας τα σωστά στον κώδικα που τα αναμένει.
    \item \textbf{Εξυπηρέτηση Αιτημάτων Κλοπής (Task Stealing Requests):} Η πιο σύνθετη λειτουργία του είναι ο χειρισμός των αιτημάτων κλοπής που καταφθάνουν από άλλες, αδρανείς διεργασίες MPI που αναζητούν φόρτο εργασίας.
\end{enumerate}

Για να διασφαλιστεί ότι η επικοινωνία στο δίκτυο δεν θα μπλοκάρει την εκτέλεση του υπόλοιπου κώδικα Python (και κυρίως για να αποφευχθεί ο εγκλωβισμός του GIL της Python από το νήμα του διακομιστή), ο σχεδιασμός επιβάλλει ρητά τη χρήση μη-μπλοκαρισμένων (non-blocking) κλήσεων του MPI. Κατά την ανάπτυξη και τη βελτιστοποίηση του torcpy, η ομάδα διαπίστωσε μέσα από τη χρήση εργαλείων αξιολόγησης από το πακέτο mpi4py.futures (όπως οι ρουτίνες run\_julia.py και perf\_primes.py) ότι οι μπλοκαρισμένες κλήσεις MPI\_Recv δημιουργούσαν τεράστια συμφόρηση. Η λύση που υιοθετήθηκε στηρίχθηκε στην ενσωμάτωση non-blocking κλήσεων σε συνδυασμό με έναν μηχανισμό συνειδητής εκχώρησης του επεξεργαστή (processor yielding). Όταν το νήμα του διακομιστή δεν έχει μηνύματα να διαχειριστεί, "κοιμάται" εθελοντικά για ελάχιστο χρονικό διάστημα (π.χ. 0.01 δευτερόλεπτα), αποδεσμεύοντας πλήρως τον πυρήνα και επιτρέποντας στα νήματα των εργατών να προχωρήσουν απρόσκοπτα με την εκτέλεση του κώδικά τους. Παρά αυτή την τεχνική, παραμένει μια εγγενής πραγματικότητα: όσο αυξάνονται τα ενεργά threads που επιχειρούν να εκτελέσουν bytecode, η απόδοση θα επηρεάζεται αναπόφευκτα από το GIL \cite{hadjidoukas2020torcpy}.

\section{Μηχανισμός Εξισορρόπησης Φόρτου}

